\begin{longtable}{|p{0.8cm}|p{3.0cm}|p{3.0cm}|p{7.0cm}|}
\caption{Etapas Metodológicas, Resultados e Mapeamento DSR}
\label{tab:etapas_metodologicas}\\
\hline
\textbf{ID} & \textbf{Atividade} & \textbf{Fase DSR} & \textbf{Resultado (Artefato) e Descrição} \\
\hline
\endfirsthead

\multicolumn{4}{c}%
{{\tablename\ \thetable{} -- continuação}} \\
\hline
\textbf{ID} & \textbf{Atividade} & \textbf{Fase DSR} & \textbf{Resultado (Artefato) e Descrição} \\
\hline
\endhead

\hline
\multicolumn{4}{r}{{Continua na próxima página}} \\
\endfoot

\hline
\endlastfoot

\textbf{M1} &
\textbf{Revisão Rápida da Literatura} &
Identificação do Problema e Objetivos &
\textbf{Catálogo de Ferramentas \acs{open_source}:} Base teórica de tecnologias e \textbf{Critérios de Avaliação} de ETL \cite[]{boulahia2024} para guiar a seleção (Capítulo 3). \\
\hline

\textbf{M2} &
\textbf{Mapeamento e Projeto da Arquitetura} &
Design e Desenvolvimento &
\textbf{Arquitetura Conceitual e Lógica:} Definição formal do \textit{stack} tecnológica (ferramentas específicas) e do fluxo de dados que será implementado no protótipo. \\
\hline

\textbf{M3} &
\textbf{Construção do Protótipo Funcional} &
Design e Desenvolvimento &
\textbf{Protótipo do Pipeline Funcionando:} Código implementado para Extração (\textit{web scraping}), scripts de Transformação (limpeza e enriquecimento) e Carga dos dados (Capítulo 4). \\
\hline

\textbf{M4} &
\textbf{Demonstração e Coleta de Métricas} &
Demonstração &
\textbf{Dados Operacionais Brutos:} Logs de \textbf{Tempo Médio de Processamento} e relatórios de \textbf{Consumo de Recursos Computacionais} (CPU/Memória) sob carga, essenciais para a avaliação. \\
\hline

\textbf{M5} &
\textbf{Avaliação Crítica e Análise Comparativa} &
Avaliação &
\textbf{Análise Crítica e Lições Aprendidas:} Documentação dos \textit{benchmarks} de desempenho, limitações de escalabilidade e a confirmação da viabilidade econômica das soluções \acs{open_source} (Capítulo 5). \\
\hline

\textbf{M6} &
\textbf{Comunicação dos Resultados} &
Comunicação &
\textbf{Trabalho de Conclusão de Curso (TCC):} Documento final que formaliza o artefato criado e o novo conhecimento gerado. \\
\hline

\end{longtable}